\documentclass{article}
\usepackage{hyperref}

\title{Keil Setup FAQ}
\author{}
\date{}

\begin{document}

\maketitle

\section*{Introduction}

This document outlines common issues and solutions encountered while setting up a project in the Keil IDE, specifically when integrating STM32CubeMX and dealing with compiler and debugger configurations.

\section*{Q: What should I do if I encounter errors related to ARM Compiler version?}

A common error is: \textit{“Target 'Project' uses ARM-Compiler 'Default Compiler Version 5' which is not available.”} This occurs because STM32CubeMX defaults to generating code for ARM Compiler Version 5, while Keil may expect Version 6. 

\begin{itemize}
    \item \textbf{Solution:} In the project options, ensure that ARM Compiler Version 6 is selected. You can adjust this under \texttt{Options for Target -> Target Tab}.
    \item If needed, reinstall or update ARM Compiler packs using the Pack Installer, ensuring both Version 5 and 6 are available.
\end{itemize}

\section*{Q: What do I do if my project cannot find necessary DLLs? (e.g., SarmCM3.dll, STLinkUSBDriver.dll)}

\begin{itemize}
    \item \textbf{SarmCM3.dll Error:} If Keil cannot locate \texttt{SarmCM3.dll} or similar files, ensure that the full path to these files is specified in \texttt{Options for Target -> Debug Tab}. If not, add it to the PATH environment variable.
    \item \textbf{STLinkUSBDriver.dll Error:} For ST-Link debugging errors related to missing DLLs (e.g., \texttt{ST-LINKIII-KEIL\_SWO.dll}), ensure that the ST-Link USB drivers are installed, and the path to \texttt{STLinkUSBDriver.dll} is correctly set.
\end{itemize}

\section*{Q: How do I solve the error “Registered ARM Compiler version not found in path” even though it builds?}

\begin{itemize}
    \item This can be resolved by ensuring that the correct path to the ARM Compiler is set in the environment variable \texttt{PATH} or configured in Keil settings. The compiler path should point to \texttt{C:\textbackslash{}Keil\textbackslash{}ARM\textbackslash{}BIN}.
    \item In \texttt{Options for Target -> Target Tab}, ensure that the correct compiler path is listed. ARM Compiler 6.22 or ARMCLANG is typically expected.
\end{itemize}

\section*{Q: What steps should I follow if Keil can't find \texttt{ST-LINKIII-KEIL\_SWO.dll}?}

\begin{itemize}
    \item Ensure that you have the ST-Link drivers installed. These can be found under the Keil installation folder, typically at \texttt{C:\textbackslash{}Keil\textbackslash{}ARM\textbackslash{}STLink}.
    \item Make sure the correct path is referenced in \texttt{Options for Target -> Debug Tab}. You can add the full path to the \texttt{ST-LINKIII-KEIL\_SWO.dll}.
    \item If ST-Link is not detected, switch to ULINK/ME for debugging, which sometimes resolves this issue.
\end{itemize}

\section*{Q: Why am I getting “Additional software components required” errors in Keil RTE?}

This usually happens when using STM32CubeMX to generate code and Keil expects other components to be available.
\begin{itemize}
    \item \textbf{Solution:} Use the Keil Pack Installer to ensure that all required packs are installed. The error message typically indicates what is missing. 
    \item Check the \texttt{Manage Run-Time Environment} (RTE) for missing components, and use the \texttt{Pack Installer} to download them.
\end{itemize}

\section*{Q: What should I select in the RTE setup for a basic STM32CubeMX project?}

Here are typical selections for a basic project:

\begin{itemize}
    \item \textbf{CMSIS Core:} Include the CMSIS core components.
    \item \textbf{Device Startup:} Choose the appropriate startup files for your STM32 device.
    \item \textbf{STM32Cube Framework:} For a HAL-based project, select STM32Cube HAL for peripheral initialization.
    \item \textbf{Debugger:} Choose ST-Link for debugging (unless you’re using another debugger like ULINK).
\end{itemize}

\section*{Q: How do I make semihosting work with Keil?}

Keil supports semihosting, but ensure the following:
\begin{itemize}
    \item Semihosting must be enabled via the project options (select \texttt{use semihosting}).
    \item Use appropriate \texttt{\_write()} and \texttt{\_read()} implementations for redirecting the standard output to the debugger's console.
\end{itemize}

\section*{Q: How do I add STM32CubeMX code generation to Keil projects?}

If STM32CubeMX fails to integrate properly:
\begin{itemize}
    \item First, generate the code using STM32CubeMX, then import it manually into Keil by copying the files over.
    \item Ensure that the include paths and all project settings match the Keil project environment.
\end{itemize}

\section*{Q: How do I reset Keil if it becomes unstable or can't find files?}

\begin{itemize}
    \item Resetting your Keil installation and reinstalling the packs can often resolve issues with missing files or paths.
    \item Check all environment variables and make sure the Keil paths are set correctly.
    \item If all else fails, re-install Keil, ensuring that all needed software packages are installed via the Pack Installer.
\end{itemize}

\section*{Q: What if Keil can't save files or throws “Member not found” errors?}

This might be due to a corrupted session or unsaved changes. If Keil locks up and throws errors such as \texttt{HRESULT DISP\_E\_MEMBERNOTFOUND}, try the following:
\begin{itemize}
    \item Close Keil and re-open the project.
    \item If the issue persists, reset your project or open a backup.
\end{itemize}

\section*{Conclusion}

These are the main issues and resolutions when setting up a Keil project for STM32 microcontrollers. With the right settings and configurations, Keil can be a powerful and stable tool for embedded development.

\end{document}
